\textbf{Paschenova křivka} vyjadřuje závislost zápalného napětí na součinu tlaku plynu $p$ a vzdálenosti elektrod $d$:
\[U_z=f(pd).\]
Pro daný plyn a materiál elektrod odpovídá stejnému $pd$ stejné $U_z$. Křivka má minimum; při menším i větším $pd$ zápalné napětí roste.

Všechny její body splňují \textbf{mezní Townsendovu podmínku samostatného výboje} $\gamma(e^{\alpha d}-1)=1$. Křivka odděluje oblast nesamostatného a samostatného výboje.

Slouží k určení zápalného napětí pro daný tlak a mezeru, případně k volbě tlaku či vzdálenosti elektrod pro požadované zápalné napětí.
